% ============================================================================
% SUBSTRATE: RUNTIME ARCHITECTURE
% ============================================================================
\subsection{Computational Substrate}
\label{sec:substrate}

Values do not execute in isolation.  The runtime organizes them into
concurrent pools, routes boolean operations to available hardware, and
manages the flow of frames through a buffered streaming pipeline.

\paragraph{Regions.}
\label{sec:regions}
A \emph{Region} is a concurrent mixing pool that accepts and returns
Value-sized frames.  Each Region maintains three tiers of storage: a
primary mixer channel (capacity 64 frames), a bounded spill channel
(capacity 256 frames), and an unbounded overflow queue.  Writes always
succeed---a frame is deposited into the primary channel if space is
available, otherwise into the spill channel, and finally into the
overflow queue.  Reads drain the tiers in the same order, returning
frames non-deterministically.

This design guarantees that no frame is ever dropped or blocked during
concurrent access.  The non-deterministic read order means that frames
deposited by different producers may interact in varying sequences,
creating a form of stochastic mixing that enables combinatorial
exploration of frame interactions without explicit scheduling.

\paragraph{Multi-Substrate Backend.}
\label{sec:multi_backend}
A \emph{Backend} abstracts the heterogeneous hardware available on the
host.  At initialization, the Backend probes for NVIDIA CUDA devices,
Apple Metal devices, and the always-present CPU, registering each as a
compute substrate.  All substrates implement a common interface with two
operations: \texttt{UniversalBitwise}, which reads a four-bit opcode
from a Value's current instruction, executes the corresponding boolean
truth table across the frame's data lanes, and advances the program
counter; and \texttt{Schedule}, which enqueues a job for asynchronous
execution on the substrate's worker path.

The Backend selects a substrate for each operation through round-robin
rotation over the registered hardware.  An exception is made for frames
that arrive over wide-area network (WAN) transport: these frames are
pinned to the substrate that first processes them, avoiding repeated
migration between devices.

\paragraph{Machine.}
The \emph{Machine} is the top-level runtime orchestrator.  It binds
together a worker pool, a transport stream, a set of Regions, and an
optional dataset source.  The worker pool allocates one goroutine per
available CPU core (with a minimum of one), and each worker draws jobs
from a shared channel.  Workers are recycled through a pool to minimize
allocation overhead, and each is wrapped in a panic-recovery guard to
prevent a single faulting job from terminating the runtime.

\paragraph{Stream and Framing.}
The \emph{Stream} provides a ring-buffered framing layer between the
Machine and its Regions.  Incoming data is written in frame-aligned
units of 1024 bytes.  Each complete frame is first committed to the ring
buffer and then fanned out to every Region, split into two configurable
sides.  Reads from the Stream return one complete frame at a time.  If
the final write is shorter than a full frame, the Stream pads the
remainder with zero bytes before committing.

\paragraph{Network Transport.}
A unified connection type, \emph{UniConn}, abstracts three transport
mechanisms: inter-process communication (IPC) via Unix domain sockets,
UDP multicast for local-network broadcast, and QUIC for reliable
wide-area transport.  All three implement the standard
\texttt{io.ReadWriteCloser} interface, allowing Values to flow across
process and network boundaries using the same read/write operations
used within a single Machine.  Frames arriving over QUIC are tagged
with an ingress marker so that the Backend can apply the pinning policy
described above.
